\documentclass[11pt]{article}
\usepackage[a4paper,margin=1in]{geometry}
\usepackage{amsmath}
\usepackage{amssymb}
\usepackage{bm}
\usepackage{mathtools}
\usepackage{booktabs}
\usepackage{hyperref}

\title{An Introduction to Krotov Optimization and Its Adaptation to Quantum Machine Learning}
\author{Codex note}
\date{\today}

\begin{document}
\maketitle

\section*{Purpose of this note}
This note is meant as a compact mathematical introduction to the Krotov optimization method and, more specifically, to the variant used in this project for quantum machine learning (QML). 

\section*{1. The basic optimal-control problem}
Krotov's method originates in optimal control. The standard setting is:
\begin{itemize}
\item a state variable $\psi(t)$ evolving over time;
\item a control $\epsilon(t)$ that influences the dynamics;
\item a cost functional $J[\psi,\epsilon]$ that should be minimized.
\end{itemize}

In a quantum setting the state often obeys a Schr\"odinger equation
\begin{equation}
i \frac{d}{dt}\psi(t) = H(\epsilon(t), t)\psi(t),
\end{equation}
with some initial condition $\psi(0)=\psi_0$. The control $\epsilon(t)$ enters through the Hamiltonian.

A typical objective has the form
\begin{equation}
J[\psi,\epsilon]
=
\Phi(\psi(T))
\;+\;
\int_0^T g_a(\epsilon(t),t)\,dt
\;+\;
\int_0^T g_b(\psi(t),t)\,dt.
\label{eq:generic_cost}
\end{equation}
Here:
\begin{itemize}
\item $\Phi(\psi(T))$ is the terminal objective, for example ``make the final state close to a target state'';
\item $g_a$ penalizes large or irregular controls;
\item $g_b$ can penalize undesirable behavior along the trajectory.
\end{itemize}

The essential point is that the optimization variable is not the state itself. The state is constrained by the dynamics. The true degree of freedom is the control $\epsilon(t)$.

\section*{2. What makes Krotov's method special}
At a high level, Krotov's method is an iterative control-update method. Starting from a current control $\epsilon^{(i)}(t)$, it constructs an improved control $\epsilon^{(i+1)}(t)$ by using:
\begin{itemize}
\item a forward propagation of the state;
\item a backward propagation of an adjoint variable, usually called the \emph{co-state} or \emph{costate};
\item an update formula that is local in time.
\end{itemize}

The reason the method is attractive is that, under appropriate assumptions, it is designed to yield \emph{monotonic improvement}: each iteration should not increase the cost functional. That is the central theoretical idea behind the method.

This is different in flavor from ordinary gradient descent. A pure gradient method says ``move in the negative gradient direction.'' Krotov's method instead derives an update rule from the structure of the constrained variational problem, combining state propagation, adjoint propagation, and a carefully chosen quadratic control penalty.

\section*{3. The forward state and the backward co-state}
The forward state $\psi(t)$ represents the physical system evolving under the current control. The backward co-state $\chi(t)$ carries sensitivity information from the objective back through the dynamics.

In the simplest terminal-cost setting, the final condition for the co-state is determined by the derivative of the terminal objective:
\begin{equation}
\chi(T) = - \frac{\partial \Phi}{\partial \langle \psi(T) |}.
\label{eq:costate_terminal}
\end{equation}
Depending on sign conventions one may write the same relation with the opposite sign. The important fact is not the sign convention itself, but the role of $\chi(T)$: it encodes how the final-time objective changes when the final state changes.

Once this terminal condition is known, the co-state is propagated backward with the adjoint dynamics. In a simple Schr\"odinger setting that means
\begin{equation}
i \frac{d}{dt}\chi(t) = H(\epsilon^{(i)}(t), t)\chi(t)
\end{equation}
or, equivalently, propagation with the adjoint propagator. Again, exact signs and adjoints depend on the chosen convention, but the structure is always the same: the co-state transports final-time sensitivity information backward to earlier times.

\section*{4. A typical Krotov update formula}
For a Hamiltonian that depends smoothly on the control, one common Krotov update has the schematic form
\begin{equation}
\epsilon^{(i+1)}(t)
=
\epsilon^{(i)}(t)
\;+\;
\frac{S(t)}{\lambda_a}
\,
\operatorname{Im}
\left\langle
\chi^{(i)}(t)
\middle|
\frac{\partial H}{\partial \epsilon}
\middle|
\psi^{(i+1)}(t)
\right\rangle .
\label{eq:krotov_update}
\end{equation}
Here:
\begin{itemize}
\item $S(t)$ is a shape function, often used to suppress updates near the time boundaries;
\item $\lambda_a$ is a step-size-like regularization parameter;
\item $\chi^{(i)}(t)$ is the co-state from the old iterate;
\item $\psi^{(i+1)}(t)$ is the forward-propagated state under the new iterate.
\end{itemize}

Two features are worth emphasizing.

First, the update is \emph{implicit-looking}: the new control appears together with the new forward state. This is one reason Krotov's method differs conceptually from a standard explicit gradient step.

Second, the update is \emph{local in time}: at each time $t$, one updates the control using state and co-state information available at that same time. In continuous control problems that is highly attractive, because it naturally respects the dynamical structure of the system.

\section*{5. Why Krotov's method fits quantum control well}
The method is especially natural in coherent quantum control because the ingredients match the physics:
\begin{itemize}
\item the forward state is the actual quantum state trajectory;
\item the co-state is another state-like object propagated with adjoint dynamics;
\item the control enters through the Hamiltonian;
\item the final objective is often expressed directly in terms of overlaps or observables.
\end{itemize}

So the control update can be written entirely in terms of inner products of propagated quantum states and generators of the controlled dynamics.

\section*{6. From continuous-time control to QML training}
The QML setting in this project is not a continuous-time pulse-shaping problem. Instead, the trainable model is a parameterized quantum circuit:
\begin{equation}
U(\bm{\theta},x)
=
U_M(\theta_M,x)\cdots U_2(\theta_2,x)U_1(\theta_1,x),
\end{equation}
where:
\begin{itemize}
\item $\bm{\theta}=(\theta_1,\dots,\theta_M)$ is the trainable parameter vector;
\item $x$ is the input data point;
\item each $U_k$ is a gate, often of the form
\begin{equation}
U_k(\theta_k,x)=\exp\!\left(-\frac{i}{2}\alpha_k(x,\theta_k)G_k\right)
\end{equation}
for some Hermitian generator $G_k$.
\end{itemize}

Given a fixed initial state $\psi_0$, the circuit produces a final state
\begin{equation}
\psi_M(x;\bm{\theta}) = U(\bm{\theta},x)\psi_0.
\end{equation}

In classification, the model prediction is then obtained from a measurement. In the simplest case one measures a Hermitian observable $O$ and maps its expectation value to a probability:
\begin{equation}
z(x;\bm{\theta}) = \langle \psi_M | O | \psi_M \rangle,
\qquad
p(x;\bm{\theta}) = \frac{z(x;\bm{\theta})+1}{2}.
\label{eq:simple_readout}
\end{equation}

The empirical loss over a dataset $\{(x_n,y_n)\}_{n=1}^N$ is usually
\begin{equation}
\mathcal{L}(\bm{\theta})
=
\frac{1}{N}\sum_{n=1}^N
\ell\bigl(p(x_n;\bm{\theta}), y_n\bigr),
\end{equation}
with $\ell$ the binary cross-entropy in this project.

\section*{7. The key adaptation: replace continuous time by gate index}
The essential conceptual move is simple:
\begin{itemize}
\item in continuous Krotov, the update is local in \emph{time};
\item in circuit QML, the natural discrete analogue is locality in \emph{gate index}.
\end{itemize}

So instead of a trajectory $\psi(t)$ we work with a sequence of forward states
\begin{equation}
\psi_0,\psi_1,\dots,\psi_M,
\qquad
\psi_{k+1}=U_{k+1}(\theta_{k+1},x)\psi_k.
\end{equation}

And instead of a continuous-time co-state $\chi(t)$ we work with a backward sequence
\begin{equation}
\chi_M,\chi_{M-1},\dots,\chi_0,
\qquad
\chi_k = U_{k+1}^\dagger(\theta_{k+1},x)\chi_{k+1}.
\end{equation}

This is the finite-dimensional discrete analogue of forward and backward propagation.

\section*{8. Terminal co-state for QML classification}
The terminal co-state is derived from the derivative of the loss with respect to the final state. For the simple readout \eqref{eq:simple_readout} and binary cross-entropy loss,
\begin{equation}
\ell(p,y) = -y\log p - (1-y)\log(1-p),
\end{equation}
one obtains
\begin{equation}
\chi_M
=
\frac{1}{2}
\frac{\partial \ell}{\partial p}
\,
O\psi_M.
\label{eq:terminal_costate_qml}
\end{equation}

This formula is important because it shows the direct analogy with continuous quantum control: the terminal co-state is still ``loss derivative times observable acting on the state.'' The only difference is that the terminal objective is now a machine-learning loss rather than a state-transfer infidelity.

If the readout contains an additional classical head, one first differentiates through that head. For example, if
\begin{equation}
m_j = \langle \psi_M | O_j | \psi_M \rangle,
\qquad
s = \bm{w}^\top \bm{m} + b,
\qquad
p=\sigma(s),
\end{equation}
then binary cross-entropy gives
\begin{equation}
\frac{\partial \ell}{\partial s}=p-y,
\end{equation}
so the terminal co-state becomes
\begin{equation}
\chi_M
=
(p-y)
\left(
\sum_j w_j O_j
\right)\psi_M.
\label{eq:terminal_costate_hybrid_head}
\end{equation}

That is exactly the form used for the Simonetti-style hybrid models in this project.

\section*{9. Gate-wise update direction}
Suppose gate $k$ depends on a scalar parameter $\theta_k$ and is generated by a Hermitian operator $G_k$:
\begin{equation}
U_k(\theta_k)=\exp\!\left(-\frac{i}{2}\theta_k G_k\right).
\end{equation}
Then
\begin{equation}
\frac{\partial U_k}{\partial \theta_k}
=
\left(-\frac{i}{2}G_k\right)U_k.
\end{equation}

In the discrete Krotov-style update used here, the contribution associated with parameter $\theta_k$ is
\begin{equation}
\Delta_k
:=
2\,\operatorname{Re}
\left\langle
\chi_k^{\text{after}}
\middle|
\left(-\frac{i}{2}G_k\right)
\middle|
\psi_k^{\text{after}}
\right\rangle ,
\label{eq:discrete_contribution}
\end{equation}
where $\psi_k^{\text{after}}$ is the forward state after the gate and $\chi_k^{\text{after}}$ is the backward state at the same location. In the implementation, this appears as a gate-local inner product between a propagated co-state and a gate-derivative generator acting on the propagated forward state.

The update then has the schematic form
\begin{equation}
\theta_k^{\text{new}}
=
\theta_k^{\text{old}}
\;-\;
\eta \, \Delta_k,
\end{equation}
up to sign convention, depending on whether one writes $\Delta_k$ as a descent direction or as a gradient-like quantity. In the code base, the quantity in \eqref{eq:discrete_contribution} is treated as the contribution to be subtracted after multiplication by the chosen step size.

\section*{10. Why this is only ``Krotov-inspired'' in QML}
It is important to be precise here. The adaptation used in QML is not a literal transcription of the original continuous-time monotonic Krotov theorem. Several things have changed:
\begin{itemize}
\item time has become a finite gate index;
\item the objective is an empirical loss over a dataset, not a single trajectory objective;
\item the updates are applied to circuit parameters, not directly to a smooth control field;
\item practical training introduces mini-batch, full-batch, or sample-wise averaging choices.
\end{itemize}

So the correct description is that this is a \emph{Krotov-inspired discrete adjoint update} for QML, not a fully general proof-carrying continuous-control Krotov scheme. Nevertheless, the structural analogy is real and useful: forward states, backward co-states, gate-local update directions, and a natural decomposition by propagation steps.

\section*{11. How the project uses the method in practice}
In this project three closely related variants are used.

\subsection*{11.1 Online Krotov}
For one training sample $(x_n,y_n)$:
\begin{enumerate}
\item forward-propagate through the gate sequence and store intermediate states;
\item construct the terminal co-state from the sample loss;
\item backward-propagate the co-state through the circuit;
\item update parameters gate by gate using the local contribution formula.
\end{enumerate}

This preserves the strongest intuition from the original Krotov idea: sequential local updates informed by forward and backward propagation. It is often very aggressive early in training, because it reacts strongly to individual samples.

\subsection*{11.2 Batch Krotov}
Instead of updating after each sample, one computes the gate-wise contribution for each sample and averages:
\begin{equation}
\bar{\Delta}_k
=
\frac{1}{N}\sum_{n=1}^N \Delta_k^{(n)}.
\end{equation}
Then one performs a single parameter update with this averaged direction.

This is closer in spirit to full-batch optimization and aligns more directly with the empirical mean loss. It is typically more stable than the online rule, but less aggressive in the very first iterations.

\subsection*{11.3 Hybrid Krotov}
The hybrid scheme used in the report combines the two phases:
\begin{itemize}
\item an early online phase to exploit the fast initial progress of sample-wise updates;
\item a later batch phase to refine the solution more smoothly.
\end{itemize}

Empirically, this matches the observed behavior of the method quite well: online Krotov is often excellent at making a strong early move into a good basin, while batch Krotov is better at reducing oscillations once the model is already close to a good classifier.

\section*{12. How non-gate parameters are handled}
Some QML models are hybrid in a literal sense: a quantum circuit produces several expectation values, and a classical affine or logistic head maps those values to a final prediction. Such classical head parameters are not themselves quantum gates, so a gate-local Krotov update does not apply directly.

In the present project this issue is handled in the cleanest possible way:
\begin{itemize}
\item gate-supported parameters receive the Krotov-style update described above;
\item purely classical non-gate parameters receive their exact loss gradient update within the same training epoch.
\end{itemize}

This is an important practical adjustment. It allows full hybrid models, such as the Simonetti classifier, to be trained fairly against Adam and L-BFGS-B without freezing part of the model merely to fit the optimizer.

\section*{13. Why this adaptation is attractive for QML}
There are several reasons the method is interesting in QML.
\begin{enumerate}
\item It respects circuit structure. The update is built from the actual gate sequence and its propagated states, not only from a black-box gradient oracle.
\item It naturally exploits generator information. If a gate has a known Hermitian generator, the local update direction is easy to express.
\item It provides an alternative to parameter-shift-based optimization. Parameter-shift computes gradients by repeated shifted evaluations; the Krotov-style update instead uses stored forward states and backward co-states.
\item It suggests architecture-aware optimization. Models with a strongly gate-centric structure may be especially suitable for this type of update.
\end{enumerate}

\section*{14. Limitations and caveats}
At the same time, several caveats are important.
\begin{itemize}
\item The classical monotonic-convergence guarantees of continuous Krotov do not transfer automatically to the discrete QML setting.
\item The method is most natural when one can reconstruct gate-local generators and intermediate states. It is therefore less black-box than generic optimizers.
\item For hybrid quantum-classical models, one must decide how to combine gate-local updates with ordinary classical gradients.
\item Empirical performance is architecture-dependent. In the experiments of this project, hybrid Krotov is clearly useful on some models and not dominant on others.
\end{itemize}

\section*{15. The short conceptual summary}
If one compresses the whole story into one sentence, it is this:

\begin{quote}
Krotov's method is an optimal-control update rule that improves a control by combining a forward trajectory, a backward sensitivity trajectory, and a local update formula; in QML, the same idea is transferred from continuous time to discrete gate position, yielding a gate-wise training rule based on forward circuit states and backward co-states.
\end{quote}

That is the right conceptual picture to keep in mind. The original method belongs to optimal control. The QML version used here is a discrete, circuit-level adaptation of the same adjoint-based logic.

\section*{16. Summary table}
\begin{table}[h]
\centering
\begin{tabular}{@{}lll@{}}
\toprule
Continuous Krotov idea & QML analogue in this project & Interpretation \\
\midrule
Time $t$ & Gate index $k$ & locality of the update \\
Control field $\epsilon(t)$ & Parameter vector $\bm{\theta}$ & optimization variable \\
State trajectory $\psi(t)$ & Forward gate states $\psi_k$ & propagated model state \\
Co-state $\chi(t)$ & Backward gate states $\chi_k$ & propagated sensitivity \\
Terminal objective $\Phi(\psi(T))$ & Sample loss $\ell(p(x),y)$ & training objective \\
$\partial H/\partial \epsilon$ & Gate generator derivative & local update operator \\
Monotonic control update & Krotov-inspired gate update & structured training step \\
\bottomrule
\end{tabular}
\caption{Dictionary from standard Krotov control language to the QML adaptation used in this project.}
\end{table}

\section*{Final remark}
For discussion with a mathematically trained audience, the most important message is that the method is not ``just another optimizer.'' Its natural home is constrained dynamical optimization. What makes it interesting for QML is precisely that a parameterized quantum circuit \emph{is} a dynamical object in discrete form: a sequence of unitary propagation steps. Once one sees that structural match, the forward/co-state/gate-local update construction becomes quite natural.

\end{document}
